\noindent\textit{2008 manuscript.}

\section{Ideal Gas}

\subsection*{Classical Limit}

Recall, Ch. 6, Sec. ``Classical Limit'' of Kittel and Kroemer \cite{CKittelHKroemer1980}, \textbf{an ideal gas is defined as a system of free noninteracting particles in the classical regime}.  \\
$f(\epsilon) \equiv $ average occupancy of an orbital at energy $\epsilon$ \\
$\epsilon \equiv $ energy of orbital occupied by 1 particle; not energy of system of $N$ particles \\
Fermi-Dirac and Bose-Einstein distribution $f(\epsilon) = \frac{1}{ \exp{ [ (\epsilon - \mu )/\tau ] } \pm 1 }$ \\
In order for $f(\epsilon) \ll 1$ $\forall \, $ orbitals, $\exp{ [ (\epsilon-\mu)/\tau ] } \gg 1\, \forall \, \epsilon$.  
\[
\Longrightarrow f(\epsilon) \simeq \exp{ [(\mu - \epsilon )/\tau]} = \lambda \exp{ (-\epsilon/\tau)} \quad \, \lambda \equiv \exp{ \left( \frac{ \mu }{\tau} \right) }
\]
$f(\epsilon)$, average occupancy of orbital of energy $\epsilon$, is classical distribution function.  

In Ch. 3 of Kittel and Kroemer \cite{CKittelHKroemer1980}, $Z_1 = n_Q V$, $n_Q = \left( \frac{M\tau}{ 2\pi \hbar^2} \right)^{3/2}$.

$\lambda = e^{\mu /\tau}$
\[
N=\lambda Z_1 = \lambda n_Q V \quad \,  \lambda = \frac{N}{n_Q V} = \frac{n}{n_Q}
\]
\[
\Longrightarrow \lambda = \exp{ \left( \frac{\mu}{\tau} \right)} = \frac{n}{n_Q}
\]

\subsubsection*{Heat Capacity}

Now 
\[
Q = Q(\tau,p) = \Lambda_p dp + C_p d\tau = dU + W = dU + p dV 
\]
Let $c\in \Sigma$ s.t. $dp(\dot{c})=0$ (constant pressure).  And so 
\[
Q(\dot{c}) = 0 + C_p d\tau(\dot{c}) = \tau d\sigma (\dot{c}) = dU(\dot{c}) + pdV(\dot{c})
\]
for $c=(\tau,0)$, $\dot{c} = \frac{\partial }{ \partial \tau} \in T\Sigma$.  

\[
\Longrightarrow C_p = \tau \left( \frac{ \partial \sigma }{ \partial \tau} \right)_p = \left( \frac{ \partial U}{ \partial \tau} \right)_p + p \left( \frac{ \partial V}{ \partial \tau} \right)_p
\]
for heat capacity at constant pressure is larger than $C_V$ because additional heat must be added to perform the work needed to expand volume of gas against constant pressure.  

Now recall for enthalpy 
\[
\begin{aligned}
  & H = U+ pV \\ 
  & H = H(\sigma, p)
\end{aligned}
\]
Then
\[
dH = dU+ pdV + Vdp = Q + Vdp 
\]
Thus, for $c\in \Sigma$, $dp(\dot{c})=0$ (constant pressure).  Hence
\[
dH(\dot{c}) = \left( \frac{ \partial H}{ \partial \tau} \right)_p = Q(\dot{c})  +  0 = C_p
\]
Hence, 
\[
C_p = \left( \frac{ \partial H}{ \partial \tau } \right)_p = \left( \frac{ \partial Q}{ \partial \tau} \right)_p
\]

Now
\[
\int_c dH = H_f - H_i = \int_c Q = \int_c C_p d\tau = C_p(\tau_f - \tau_i)
\]
(if $C_p$ constant over desired temperature range)



\subsection*{Reversible Isothermal Expansion}

$Q=0$, insulated gas, no heat flow to or from the gas (``adiabatic'') \\
$\sigma$ constant in system isolated from reservoir, if expansion reverisble (slowly)

\subsubsection*{What is the pressure after expansion?}

Remember 
\[
\frac{C_p}{C_V} = \frac{ \frac{5}{2} N }{ \frac{3}{2} N } = \frac{5}{3} = \gamma 
\]
So let 
\[
\frac{1}{\gamma- 1} = \frac{3}{2}
\]
\begin{equation}
  \sigma(\tau,V) = N ( \ln{ \tau^{\frac{1}{ \gamma - 1 } } + \ln{V} + \text{ constant } } ) \quad \quad \, (61)
\end{equation}

\begin{equation}
  \ln{ \tau^{\frac{1}{ \gamma -1 } }V} = \text{ constant or } \tau^{ \frac{1}{\gamma-1}}V \text{ constant } \quad \quad \, (62)
\end{equation}

\begin{equation}
  \Longrightarrow \tau_1^{\frac{1}{\gamma-1}}V_1 = \tau_2^{\frac{1}{\gamma-1}}V_2 \quad \quad \, (63)
\end{equation}
for ideal monoatomic gas.  

Use $pV = N\tau$

\begin{equation}
\Longrightarrow \frac{ \tau_1^{\frac{\gamma}{\gamma-1}} }{p_1} = \frac{ \tau_2^{\frac{\gamma}{\gamma-1}} }{p_2} \quad \quad \, (64)
\end{equation}
or $p_1^{\frac{1}{\gamma-1}}V_1^{\frac{\gamma}{\gamma-1}} = p_2^{\frac{1}{\gamma-1}}V_2^{\frac{\gamma}{\gamma-1}}$ or $p_1V_1^{\gamma} = p_2 V_2^{\gamma}$

I will recap Problem 10 and my solution.  

\emph{Isentropic relations of ideal gas.}

\begin{enumerate}
  \item[(a)] $\gamma = \frac{C_p}{C_V}$.  For isentropic process, $pV^{\gamma} = p_i V_i^{\gamma}$.  

Then, essentially and equivalently, take the exterior derivative:
\[
V^{\gamma}dp + \gamma p V^{\gamma -1} dV = 0 \Longrightarrow \frac{dp}{p} + \frac{\gamma dV}{ V} = 0 
\]

For 
\[
\tau V^{\gamma-1} = \tau_i V_i^{\gamma-1}
\]
then taking $d$:
\[
d\tau V^{\gamma-1} + \tau(\gamma -1) V^{\gamma -2} dV = 0 \Longrightarrow \frac{d\tau}{\tau } + \frac{ \gamma -1}{V} dV = 0 
\]

For 
\[
\tau^{ \frac{ \gamma}{ 1 - \gamma }} p = \tau_i^{\frac{ \gamma}{ 1 - \gamma } } p_i
\]
then taking $d$:
\[
dp \tau^{ \frac{\gamma}{1- \gamma }} + \frac{ \gamma }{ 1 - \gamma} \tau^{ \frac{2 \gamma -1 }{ 1- \gamma }} d\tau p = 0 \Longrightarrow \frac{dp}{p} + \frac{ \gamma }{ 1 - \gamma} \frac{d\tau}{\tau} = 0 
\]
  \item[(b)] isentropic bulk moduli
\[
B_{\sigma} = -V \left( \frac{ \partial p }{ \partial V} \right)_{\sigma} = \gamma \frac{ p_i V_i^{\gamma }}{ V^{\gamma }} = \gamma p
\]
since $p = \frac{ p_i V_i^{\gamma }}{ V^{\gamma}}$
\[
\boxed{ B_{\sigma} = \gamma p }
\]
EY : 20150606 I think when one consider small, linear longitudinal perturbations of the gas system, with pressure being the external restoring force, then sound waves propagate (correct me if I'm wrong) and this is the way to derive $B_{\sigma}$.  

isothermal bulk moduli $B_{\tau} = -V \left( \frac{ \partial p }{ \partial V }\right)_{\tau} = \frac{n\tau}{V} = p$
 \\
velocity of sound in gas is $c = \left( \frac{ B_{\sigma}}{ \rho } \right)^{1/2} = \left( \frac{ \gamma p }{ \rho } \right)^{1/2}$

for ideal gas of molecules of mass $M$, $pV = N\tau$

\end{enumerate}

Manifold interpretation; I'm using Chapter 5 Applications in Physics, Section A Thermodynamics of \cite{BSchutz1980}.  Consider a 2-dimensional manifold of equilibrium states $\mathcal{M}$, s.t. $\text{dim}\mathcal{M} =2$.  Local coordinates are $(U,V)$ with $U$ being a global coordinate.  

Now $U = C_V \tau$ so use $\tau$ as a global coordinate, i.e. the local coordinates of $\mathcal{M}$ can be $(\tau,V)$.  

Consider a curve in $\mathcal{M}$ parametrized by $\lambda$: $(\tau(\lambda), V(\lambda))$.  The corresponding tangent vector $X = \dot{\tau} \frac{ \partial }{ \partial \tau } + \dot{V} \frac{ \partial }{ \partial V} \in \mathfrak{X}(\mathcal{M})$ has components in coordinate vector basis $\begin{aligned} & \quad \\
  & \dot{\tau} = \dot{\tau}(\lambda) \\
  & \dot{V} = \dot{V}(\lambda) \end{aligned}$

Recall
\[
\frac{d\tau}{\tau} + \frac{\gamma- 1}{V} dV =0 
\]
Consider applying this 1-form onto $X$:
\[
\left( \frac{d\tau}{\tau} - \frac{ \gamma - 1}{V}dV \right)(X) = \frac{ \dot{\tau}}{\tau} + \frac{\gamma -1}{ V} \dot{V} = 0 \Longrightarrow \frac{d}{d\lambda}(\tau V^{\gamma-1}) = 0 
\]

Then the equalities of the endpoints of this curve $(\tau(\lambda),V(\lambda))$ are the equalities above.  The interpretation is that the isentropic process draws out a curve in $\mathcal{M}$ and can be written as a curve or as a tangent vector field (specifically a section of $T\mathcal{M}$).  


\begin{align}
  \tau_1 V_1^{\gamma -1 } & = \tau_2 V_2^{\gamma -1 } \quad \quad \, (66) \\
  \tau_1^{ \frac{\gamma}{1-\gamma}} P_1 & = \tau_2^{\frac{\gamma}{1- \gamma}} P_2 \quad \quad \, (67) \\ 
P_1 V_1^{\gamma}   & = P_2 V_2^{\gamma} \quad \quad \, (68)
\end{align}
for, indeed, 
\[
\begin{gathered}
  p_1 V_1^{\gamma} = p_2 V_2^{\gamma} \xrightarrow{ pV=N\tau} p_1 \left( \frac{ N\tau_1}{p_1} \right)^{\gamma} = p_2 \left( \frac{ N\tau_2}{p_2} \right)^{\gamma} \text{ so } \\
  p_1^{1-\gamma} \tau_1^{\gamma} = p_2^{1-\gamma} \tau_2^{\gamma} \text{ or } p_1 \tau_1^{\frac{ \gamma}{1- \gamma} } = p_2 \tau_2^{\frac{ \gamma}{1- \gamma} } \\
\Longrightarrow  p_1^{\frac{1-\gamma}{\gamma} }\tau_1 = p_2^{\frac{1-\gamma}{\gamma} }\tau_2 \text{ or } \left( \frac{p_1}{p_2} \right)^{\frac{1-\gamma}{\gamma} } = \frac{\tau_2}{\tau_1} = \left( \frac{p_2}{p_1} \right)^{\frac{\gamma-1}{\gamma}}
\end{gathered}
\]


$T_1 = 300 \, K \, , \quad \,  V_1/V_2 = \frac{1}{2}, \, \quad $ 
\begin{equation}
T_2 = \left( \frac{1}{2} \right)^{2/3} (300 \, K) = 189 \, K \quad \quad \, (69)
\end{equation}

The gas is cooled in expansion process by 
\begin{equation}
T_1 - T_2 = 300 \, K - 189 \, K = 111 \, K \quad \quad \, (70)
\end{equation}
Expansion at constant entropy is important e.g. methods of refridgeration.  

\subsubsection*{What is the change in energy in the expansion?}

\[
U_2 - U_1 = C_V(\tau_2- \tau_1) = \frac{1}{ \gamma -1 } N(\tau_2 - \tau_1)
\]



\subsection*{Problems}

\solutionhead{1} \textbf{Derivative of Fermi-Dirac function}.  Recall $ f = \frac{1}{ \exp{ [(\epsilon - \mu )/\tau] } + 1 } = (e^{(\epsilon-\mu)/\tau} + 1 )^{-1}$  
\[
\partial_{\epsilon}f = -1 (e^{ (\epsilon - \mu )/\tau } + 1)^{-2} ( e^{ (\epsilon - \mu)/\tau } ) \left( \frac{1}{\tau} \right)
\]
\[
\partial_{\epsilon}f(\epsilon = \mu) = \frac{-1}{4\tau}
\]



\solutionhead{2} \textbf{Symmetry of filled and vacant orbitals}.   $\epsilon = \mu + \delta$
\[
f(\epsilon) = f(\mu + \delta) = \frac{1 }{ e^{\delta/\tau} + 1 } = \frac{e^{-\delta /\tau } }{ 1 + e^{-\delta/\tau} } = 1 + \frac{-1}{ e^{-\delta/\tau} + 1 }  = 1 - f(\mu - \delta)
\]





\solutionhead{3} \textbf{Distribution function for double occupancy statistics}.   \begin{itemize}
\item[(a)]
$ \eta = 1 + \lambda e^{-\epsilon/\tau} + \lambda^2 e^{-2\epsilon/\tau}$ where $\lambda = e^{\mu/\tau}$.  
\[
\langle N \rangle = \lambda \partial_{\lambda} \ln{\zeta} = \frac{\lambda}{ \zeta} ( e^{-\epsilon/\tau} + 2\lambda e^{-2\epsilon/\tau} ) = \frac{ \lambda e^{-\epsilon/\tau} + 2 \lambda^2 e^{-2\epsilon/\tau} }{\zeta}
\]
\item[(b)] $\zeta = 1 + 2 \lambda e^{-\epsilon/\tau} + \lambda^2 e^{-2\epsilon/\tau}$
\[
\langle N \rangle = \lambda \partial_{\lambda} \ln{\lambda} = \frac{ 2 \lambda e^{-\epsilon/\tau} + 2\lambda^2 e^{ - 2\epsilon/\tau} }{ \zeta} = \frac{ 2\lambda (e^{-\epsilon/\tau} ) (1 + \lambda e^{-\epsilon/\tau} ) }{ \zeta}
\]
\end{itemize}

\solutionhead{4} \textbf{Energy of gas of extreme relativistic particles}.  For $\epsilon_p \simeq p$, $\sum_s e^{-p/tau} =Z$.  

With the factor 2 for the 2 possible polarizations,
\[
\begin{aligned}
Z & = (2) \frac{4\pi }{8} \int_0^{\infty} p^2 dp e^{-p/\tau} = \pi \lbrace  \left. p^2 e^{-p/\tau} (-\tau) \right|_0^{\infty} - \int_0^{\infty} 2 p e^{-p/\tau} (-\tau) \rbrace = 2 \tau \pi \int_0^{\infty} p e^{-p/\tau} dp = \\ 
& = = 2 \tau \pi \lbrace \left. p e^{-p/\tau} (-\tau) \right|_0^{\infty} - \int_0^{\infty} e^{-p/\tau} (-\tau) dp \rbrace = 2 \tau^3 \pi 
\end{aligned}
\]
\[
U =  \tau^2 \frac{ \partial \ln{Z} }{ \partial \tau} = \tau^2 \partial_{\tau} \lbrace \ln{ (2\pi \tau^3  )} \rbrace = \tau^2 \partial_{\tau} (3\ln{\tau}) = \boxed{ U = 3\tau }
\]




\solutionhead{5} \textbf{Integration of the thermodynamic identity for an ideal gas}.  For constant $N$, recall
\[
d\sigma = \frac{dU}{\tau} + \frac{ p dV}{\tau}  = \frac{1}{\tau} \left( \frac{ \partial U}{ \partial \tau} \right)_V d\tau + \frac{1}{\tau} \left( \frac{ \partial U}{ \partial V} \right)_{\tau} dV + \frac{p dV}{\tau}
\]
$C_V = \left( \frac{ \partial U}{ \partial \tau} \right)_V$, and for an ideal gas $pV = N\tau$.  

\[
\int d\sigma = \sigma = C_V \ln{\tau} + N \ln{V} + \int \frac{1}{\tau} \left( \frac{ \partial U}{ \partial \tau} \right)_{\tau} dV + \sigma_1
\]
Now $U = \frac{3}{2} \tau$ for an ideal gas, so $\left( \frac{\partial U}{\partial V}\right)_{\tau} =0$.  
\[
\Longrightarrow \sigma = C_V \ln{\tau} + N \ln{V} + \sigma_1 
\]
$\sigma_1$ independent constant of $\tau$ and $V$.  










\solutionhead{6} \textbf{Entropy of mixing}.
\begin{correction}{empty solution}
The 2008 manuscript stops at this heading.
The mixing entropy, the identical-gas cancellation, and the ideal-mixture chemical potential are derived in \S\ref{sec:binary-mixtures}.
\end{correction}   







\solutionhead{7} \textbf{Relation of pressure and energy density}.  

\begin{itemize}
\item[(a)] Recall that $U = U(\sigma, V, N)$.  $p = - \left( \frac{ \partial U }{ \partial V} \right)_N$, and \\
$U = \frac{ \sum_s \epsilon_s e^{-\epsilon_s/\tau}}{ Z}$

\[
\frac{ \partial U}{ \partial V} = \sum_s \left( \left( \frac{ \partial \epsilon_s}{ \partial V} \right)_N e^{-\epsilon_s/\tau} + \epsilon_s \left( \frac{- 1}{ \tau} \left( \frac{ \partial \epsilon_s }{\partial V} \right)_N \right) e^{-\epsilon_s/\tau} \right) /Z - \sum_s  \left( - \frac{ \partial \epsilon_s}{ \partial V}\right) \left( \frac{1}{\tau } \right) e^{- \epsilon_s/\tau} \sum_{s_2} \epsilon_{ s_2} e^{ - \epsilon_s/\tau} /Z^2
\]
Now $p = -\left( \frac{ \partial U}{ \partial V} \right)_N$.  

So if the system is in state $s$; then $p_s = - \left( \frac{ \partial \epsilon_s}{ \partial V} \right)_N$.  
\[
\Longrightarrow \boxed{ \langle p \rangle = p = \frac{ \sum_s - \left( \frac{  \partial \epsilon_s }{ \partial V} \right)_N e^{-\epsilon_s/\tau}}{ Z } }
\]

\item[(b)] Now $\epsilon_s = \frac{ \pi}{ L^2} \frac{ n^2 \hbar^2}{ 2 M} = \frac{ \pi}{ V^{2/3} } \frac{ n^2 \hbar^2 }{ 2 M} $

\[
\left( \frac{ \partial \epsilon_s}{ \partial V }\right)_N = - \frac{2}{3} V^{-5/3} \frac{ \pi n^2 \hbar^2}{ 2M} = \boxed{ -\frac{2}{3V} \epsilon_s }
\]
\item[(c)] 
\[
p = \frac{ \sum_s \frac{-2 \epsilon_s}{3V } e^{-\epsilon_s /\tau } }{ Z} = \boxed{ \frac{2}{3} \left( \frac{1}{V} \right) U }
\]
\end{itemize}

\solutionhead{8} \textbf{ Time for a large fluctuation}.  

\begin{itemize}
\item[(a)] Recall $\sigma = N [ \ln{ \left( \frac{ n_Q}{ n} \right)  + \frac{5}{2}  }]$.  $n_Q = \left( \frac{ M \tau}{2\pi \hbar^2} \right)^{3/2}$.  

$\sigma_f = N [ \ln{ \left( \frac{ n_{QA} 2 V }{ N} \right) + \frac{5}{2} } ]$ \\
$\sigma_i = N [ \ln{ \left( \frac{n_{QA} V}{N} \right) + \frac{5}{2} } ]$

Now
\[
g \sim e^{\sigma} \sim \exp{ \left( \ln{ \left( \frac{n_Q}{n} \right)^N} + \frac{5N}{2} \right)} = e^{\frac{5N}{2}} \left( \frac{N_Q}{n} \right)^N
\]
\[
n_Q = 4 \times \frac{ (938 \, MeV/c^2)( (0.8617 \times 10^{-4} \, eV)/K)(300 \, K) }{ 2 \pi (6.582122 \times 10^{-22} \, MeV \cdot s )^2 } \left( \frac{ 1 \, MeV}{ 10^6 \, eV } \right) \left( \frac{ 1 \, c }{ 3\times 10^8 \, m/s } \right)^2 = 7.88 \times 10^{30}/m^3
\]

Now $PV = N\tau$.  Then
\[
\frac{ P}{\tau} = \frac{N}{V} = \frac{ ( 1 \, atm )( \frac{ 1.013 \times 10^5 \, N/m^2 }{ 1 \, atm} ) \left( \frac{ kg \cdot m/s^2 }{ 1 \, N } \right) }{ 1.381 \times 10^{-23} \, J/K  (300 \, K ) } = 2.445 \times 10^{25} / m^3
\]
Now $1 \, L = 10^{-3} \, m^3$ and so for $0.1 \, L$, 
\[
(2.445 \times 10^{25} \, \frac{1}{m^3} ) ( 10^{-4}\, m^3 ) = 2.445 \times 10^{21} 
\]
With 
\[
\left( \frac{ 7.88 \times 10^{30} / m^3}{ 2.445 \times 10^{25} \, /m^3 } \right) = 3.22 \times 10^5
\]
\[
\Longrightarrow g \sim e^{ \frac{5}{2} ( 2.445 \times 10^{21} ) } (3.22 \times 10^5 )^{2.445 \times 10^{21} }
\]
\item[(b)]
\item[(c)]
\end{itemize}






\solutionhead{9} \textbf{Gas of atoms with internal degree of freedom}.  

For an ideal monatomic gas, assume noninteracting.  
\begin{itemize}
\item[(a)] $\lambda_{int} = \lambda$.   $\lambda = \exp{(\mu /\tau)}$ ideal gas.  $\mu = \tau \ln{ ( n/n_Q)}$; \quad \, $n_Q = \left( \frac{ M\tau}{ 2 \pi \hbar^2} \right)^{3/2}$.  

\[
\boxed{ \lambda_{ext} = \exp{ (-\Delta/\tau)} \text{ or } 1  }
\]
$Z_1$ is the usual canonical partition function, $Z_1 = n_Q V$ where $n_Q = \left( \frac{M\tau}{ 2 \pi \hbar^2 } \right)^{3/2}$  
\[
\begin{gathered}
  \mathcal{Z}_1 = \sum_s ( \lambda \exp{ ( - \epsilon_s/\tau) } (1) + \lambda \exp{ \left( \frac{-\epsilon_s}{\tau} \right)} \exp{ \left( \frac{ - \Delta }{\tau} \right) } )Z_1  = \lambda ( 1 + e^{-\Delta/\tau } ) Z_1
\end{gathered}
\]

\[
( \mathcal{Z}_1 )^N = (\lambda ( 1 + e^{-\Delta/\tau} )Z_1)^N 
\]
\item[(b)]
\item[(c)]
\end{itemize}

\solutionhead{10}  \textbf{Isentropic relations of ideal gas}. 
\begin{itemize}
\item[(a)] 
Isentropic process, so $pV^{\gamma} = p_i V_i^{\gamma}$.  
\[
V^{ \gamma} dp + \gamma p V^{\gamma -1} dV = 0 \Longrightarrow \boxed{ \frac{dp}{p} + \frac{ \gamma}{ V} dV = 0 }
\]
Dealing with an ideal gas, $pV = N\tau$ still applies.  

\[
\begin{gathered}
  \tau V^{\gamma - 1 } = \tau_i V_i^{\gamma -1 } \\ 
  d\tau V^{\gamma - 1 } + \tau (\gamma - 1 ) V^{\gamma - 2 } dV = 0 \\
\Longrightarrow \boxed{   \frac{d\tau}{\tau} + \frac{ (\gamma - 1)}{ V} dV = 0  }
\end{gathered}
\]
Using $pV = N\tau$ again, note that $p^{1- \gamma} \tau^{\gamma} = \text{ constant }$.  
\[
(1- \gamma ) p^{-\gamma} dp \tau^{\gamma} + p^{1-\gamma} \gamma \tau^{\gamma - 1} d\tau =0 \Longrightarrow \boxed{ \frac{dp }{p} + \frac{ \gamma}{ 1 - \gamma} \frac{d\tau}{ \tau} = 0 }
\]
\item[(b)] Using $p = \frac{ p_i V_i^{\gamma}}{ V^{\gamma}}$, 
\[
\begin{gathered}
  \frac{ \partial p}{ \partial V} = \frac{- \gamma p_i V_i^{\gamma}}{ V^{\gamma + 1 } } \\
  \Longrightarrow -V \frac{ \partial p }{ \partial V} = \frac{ \gamma p_i V_i^{\gamma} }{V^{\gamma}} =  \gamma p
\end{gathered}
\]
So that $B_{\sigma} = -V(\partial p/ \partial V)_{\sigma} = \gamma p$, the isentropic bulk moduli.  

\[
B_{\tau} = -V \left( \frac{ \partial P }{ \partial V} \right)_{\tau} = \frac{ n \tau}{V} = p
\]
since 
\[
\begin{gathered}
  pV = n\tau \\ 
  p = \frac{ n \tau}{V} \\
  \frac{ \partial p}{\partial V} = \frac{- n \tau}{ V^2}
\end{gathered}
\]
\end{itemize}

